\section{A power graph from a growth function}\label{sec:power}

We now turn the growth function into a single-fold definition of
exponentiation.  The two sides of its growth estimate serve different
purposes.  The upper bound limits the indices of solutions to two Pell
equations, so that congruences determine those indices.  The lower bound
ensures that the intended Pell solutions satisfy the imposed limit.
The use of Pell sequences to recover exponentiation belongs to the
classical Diophantine representation method.  The paired Pell equations,
index congruences, and quotient used here adapt Matiyasevich's exponential
single-fold construction \cite[Section~4]{mati1974}, as presented in
\cite[Section~6, Lemma~6.1 and Theorem~5]{cco2024}.
Our cutoff is supplied by the complete
single-fold polynomial graph of Proposition~\ref{prop:growth}.
We begin with the elementary facts about Pell equations needed for this
argument.

\begin{lemma}\label{lem:pell-sequences}
For an integer $a\ge2$, put $D=a^2-1$ and define integers
$x_i(a),y_i(a)$ by
\[
 (a+\sqrt D)^i=x_i(a)+y_i(a)\sqrt D\qquad(i\ge0).
\]
Every solution in positive integers of $x^2-Dy^2=1$ is
$(x_i(a),y_i(a))$ for exactly one $i\ge1$.  Moreover,
\begin{equation}\label{eq:pell-bounds}
 y_i(a)\equiv i\pmod{a-1},\qquad
 (2a-1)^{i-1}\le y_i(a)\le(2a)^{i-1}\quad(i\ge1).
\end{equation}
If $b,u\ge2$ and $n\ge0$ are integers, then
\begin{equation}\label{eq:pell-ratio}
 b^n\le \frac{y_{n+1}(bu)}{y_{n+1}(u)}
 \le b^n\left(1+\frac1{2u-1}\right)^n.
\end{equation}
\end{lemma}

\begin{proof}
The integer $D$ lies strictly between $(a-1)^2$ and $a^2$, so it is
not a square.  Write $\alpha=a+\sqrt D$.  Its inverse is
$a-\sqrt D$, and both belong to $\ZZ[\sqrt D]$.
Every positive power of $\alpha$ has positive integer coefficients
and norm one.
For any positive solution, $\beta=x+y\sqrt D>1$ has inverse
$x-y\sqrt D>0$.  Choose the unique integer $i\ge0$ for which
\[
 1\le\gamma:=\beta\alpha^{-i}<\alpha.
\]
The coefficients of $\gamma=x'+y'\sqrt D$ are integers and its norm is
one.  If $\gamma>1$, then
\[
 x'=\frac{\gamma+\gamma^{-1}}2>0,\qquad
 y'=\frac{\gamma-\gamma^{-1}}{2\sqrt D}>0.
\]
Thus $y'\ge1$ and $x'^2=1+Dy'^2\ge a^2$, which implies
$\gamma\ge\alpha$, a contradiction.  Consequently $\gamma=1$ and
$\beta=\alpha^i$.  Since $\beta>1$, the index is positive; it is unique
because $\alpha>1$.

Taking coefficients gives
\[
 y_0(a)=0,\quad y_1(a)=1,\quad
 y_{i+2}(a)=2a\,y_{i+1}(a)-y_i(a).
\]
Reduction modulo $a-1$ proves the congruence in
\eqref{eq:pell-bounds} by induction.  For $i\ge1$ define
$R_i(a)=y_{i+1}(a)/y_i(a)$.  The recurrence gives
\[
 R_1(a)=2a,\qquad R_{i+1}(a)=2a-\frac1{R_i(a)}.
\]
Induction therefore yields $2a-1\le R_i(a)\le2a$.  Multiplying these
bounds proves the remaining assertions of \eqref{eq:pell-bounds}.
In particular,
\begin{equation}\label{eq:pell-index-lower}
 y_{n+1}(a)\ge(2a-1)^n\ge n+1\qquad(n\ge0).
\end{equation}

We also have $R_i(bu)\ge bR_i(u)$ for every $i\ge1$.
There is equality for $i=1$, and the induction step follows from
\[
 R_{i+1}(bu)-bR_{i+1}(u)
 =\frac b{R_i(u)}-\frac1{R_i(bu)}\ge0.
\]
On the other hand, the bounds just proved give
\[
 \frac{R_i(bu)}{R_i(u)}\le\frac{2bu}{2u-1}.
\]
Multiplication for $i=1,\ldots,n$ proves \eqref{eq:pell-ratio}.
For $n=0$ both products are empty and $y_1(bu)=y_1(u)=1$,
so the same assertion holds.
\end{proof}

The next proposition uses only the growth estimate and the uniqueness
of the entire witness for its defining system.  In particular, no
monotonicity hypothesis on the growth function is needed.

\begin{proposition}\label{prop:power}
Suppose that $g:\NN\to\NN$ has a complete single-fold graph and that
there are fixed constants $\delta>0$ and $u_0\in\NN$, with $u_0\ge4$,
such that
\begin{equation}\label{eq:power-growth-hypothesis}
 \delta(\log u)^2<\log g(u)<u\log u
 \qquad(u\ge u_0).
\end{equation}
Then the relation
\[
 \Pow(b,n,c)\quad\Longleftrightarrow\quad
 b\ge2\ \hbox{ and }\ c=b^n
 \qquad(b,n,c\in\NN)
\]
has a complete single-fold representation.  In particular, for each
$b\ge2$ and $n\in\NN$, the output $c$ and all auxiliary variables
are jointly unique.
\end{proposition}

\begin{proof}
The ratio in Lemma~\ref{lem:pell-sequences} suggests taking $c$ as the
integer quotient of $y_{n+1}(bu)$ by $y_{n+1}(u)$. To make its error
less than one without introducing a free scale, we prescribe $u$ as a
polynomial in $c,b,n$. Once the cutoff and congruences have fixed both
Pell indices at $n+1$, the lower ratio bound gives $c\ge b^n$, allowing
the factor $c+1$ to control the error. For the intended output $c=b^n$,
a sufficiently large fixed power of $c+1$ also makes $g(u)$ exceed both
intended Pell coordinates.

Choose a fixed positive integer
\begin{equation}\label{eq:power-ell}
 \ell>\frac3{\delta\log2}.
\end{equation}
Use the following finite system, in which all displayed variables are
natural numbers and $v=g(u)$ denotes one call to the given single-fold
system:
\begin{align}
 &b\ge2,\qquad
 u=(c+1)^\ell(n+1)(b+1)+u_0,\qquad v=g(u),
       \label{eq:power-input}\\
 &X^2-\bigl((bu)^2-1\bigr)F^2=1,\qquad
 Y^2-(u^2-1)H^2=1,                         \label{eq:power-pell}\\
 &0<F<v,\qquad 0<H<v,                     \label{eq:power-cutoff}\\
 &F=n+1+(bu-1)q,\qquad H=n+1+(u-1)q',      \label{eq:power-residue}\\
 &F=cH+r,\qquad r<H.                      \label{eq:power-division}
\end{align}
The exponent $\ell$ is fixed, so the formula for $u$ is polynomial.
As in Section~\ref{sec:conventions}, each inequality contributes its
unique slack variable.  The constants in
\eqref{eq:power-growth-hypothesis} determine the choice of $\ell$.
Only the fixed integers $\ell,u_0$ enter the system; $\delta$ and
logarithms occur only in its proof.

\smallskip
\noindent\emph{Every solution has the prescribed output.}
First observe the numerical bound
\begin{equation}\label{eq:pell-index-cutoff}
 u^u<(2u-1)^{u-1}\qquad(u\in\NN,\ u\ge4).
\end{equation}
Indeed, let $f(u)=(2-1/u)^{u-1}/u$.  Then $f(4)=343/256>1$, and
\[
 \frac{f(u+1)}{f(u)}
 =\frac u{u+1}
   \frac{(2-1/(u+1))^u}{(2-1/u)^{u-1}}
 >\frac{2u-1}{u+1}>1\qquad(u\ge4).
\]
This proves \eqref{eq:pell-index-cutoff}.

In a solution to \eqref{eq:power-input}--\eqref{eq:power-division},
$u\ge u_0\ge4$.  The positive coordinates $F,H$ make $X,Y$ positive
as well.  Lemma~\ref{lem:pell-sequences} therefore gives unique
indices $i,j\ge1$ with
\[
 F=y_i(bu),\qquad H=y_j(u).
\]
By \eqref{eq:power-growth-hypothesis} and
\eqref{eq:pell-index-cutoff},
\[
 F,H<v<u^u<(2u-1)^{u-1}.
\]
If either index were at least $u$, the lower bound in
\eqref{eq:pell-bounds} would contradict this strict inequality.
Hence $1\le i,j<u$.

The formula for $u$ also gives $1\le n+1<u-1$.
The congruence for $H$ in \eqref{eq:power-residue} forces
$j\equiv n+1\pmod{u-1}$.  Among $1,\ldots,u-1$ exactly the
integer $n+1$ has that residue.  Similarly,
$i\equiv n+1\pmod{bu-1}$, and both integers lie between $1$ and
$u-1<bu-1$.  Thus
\begin{equation}\label{eq:pinned-pell-indices}
 i=j=n+1.
\end{equation}
We have now determined the two Pell indices.  It remains to identify
the integer quotient imposed by \eqref{eq:power-division}.

By \eqref{eq:pell-ratio},
\[
 b^n\le F/H\le b^n(1+z)^n,
 \qquad z=\frac1{2u-1}.
\]
Since $0\le r<H$, the division equation gives
$c=\lfloor F/H\rfloor\ge b^n$.
We have $nz<1$, and the binomial theorem followed by a geometric
series gives
\[
 (1+z)^n-1
 \le\sum_{h=1}^{n}(nz)^h
 \le\frac{nz}{1-nz}.
\]
This also holds when $n=0$, with an empty sum.  It follows that
\begin{equation}\label{eq:power-quotient-error}
 0\le \frac FH-b^n
 \le\frac{b^n n}{2u-1-n}
 \le\frac{cn}{2u-1-n}<1.
\end{equation}
For the last inequality, the formula for $u$ and $\ell\ge1$ give
$u\ge(c+1)(n+1)+4$, whence
\[
 2u-1-n-cn\ge cn+2c+n+9>0.
\]
The integer quotient is consequently $c=b^n$.

\smallskip
\noindent\emph{The prescribed output has a witness.}
Fix $b\ge2$ and $n\ge0$, set $c=b^n$, and define $u$ by
\eqref{eq:power-input}.  Take the unique output and complete witness
for $v=g(u)$, and put
\[
 \begin{split}
 F=y_{n+1}(bu),&\qquad X=x_{n+1}(bu),\\
 H=y_{n+1}(u),&\qquad Y=x_{n+1}(u).
 \end{split}
\]
These coordinates satisfy the Pell equations.
The congruences in \eqref{eq:pell-bounds}, together with
\eqref{eq:pell-index-lower}, give the nonnegative integer quotients
\[
 q=\frac{F-n-1}{bu-1},\qquad
 q'=\frac{H-n-1}{u-1}.
\]
The ratio estimate and the same calculation as
\eqref{eq:power-quotient-error}, now with $c=b^n$, give
$0\le F/H-c<1$.  Thus $r=F-cH$ is a natural number less than $H$.

We must still check that these particular Pell coordinates lie below
$v$.  From $c=b^n$ and the definition of $u$ we obtain
\[
 \log u\ge\ell n\log b\ge\ell n\log2.
\]
Also $b<u$ and $2<u$, so $\log(2bu)\le3\log u$.
The upper bound in \eqref{eq:pell-bounds} now gives
\begin{equation}\label{eq:power-existence-cutoff}
 \log F\le n\log(2bu)
 \le\frac3{\ell\log2}(\log u)^2
 <\delta(\log u)^2<\log v.
\end{equation}
The same estimate applies to $H$, using
$H\le(2u)^n\le(2bu)^n$.  Hence $0<F,H<v$, as required.
When $n=0$, these statements give explicitly
\[
 c=1,\quad F=H=1,\quad X=bu,\quad Y=u,
 \quad q=q'=r=0;
\]
the strict cutoff still follows from
\eqref{eq:power-existence-cutoff}.

\smallskip
\noindent\emph{The whole witness is unique.}
For fixed $b,n$, the soundness argument first fixes $c$.
The formula in \eqref{eq:power-input} then fixes $u$, after which the
single-fold hypothesis fixes $v$ and every auxiliary variable in its
call.  Equation~\eqref{eq:pinned-pell-indices} fixes $F,H$.
Their Pell equations fix the natural coordinates $X,Y$, since a
nonnegative square root is unique.  The two residue equations fix
$q,q'$, and the division equation fixes $r$.
All inequality slacks are then fixed.  Any variables introduced to
name polynomial arguments are fixed by their defining equalities.
The indices $i,j$ used in the proof introduce no variables in the
system.  Thus the complete tuple is unique.  Finally, the constraint
$b\ge2$ excludes all solutions at $b=0,1$.
\end{proof}

Applying Proposition~\ref{prop:power} to the function supplied by
Proposition~\ref{prop:growth} gives the power graph used in the rest
of the proof.

\section{Testing containment of binary digits}\label{sec:subset}

For $x,y\in\NN$, write $x\bits y$ if every binary digit equal to one
in $x$ is also equal to one in $y$.  More explicitly, if
\[
 x=\sum_{i\ge0}x_i2^i,\qquad y=\sum_{i\ge0}y_i2^i,
 \qquad x_i,y_i\in\{0,1\},
\]
then $x\bits y$ means $x_i\le y_i$ for every $i$.
In particular, $x\bits y$ implies $x\le y$.
We will use this relation to restrict the permitted positions of
nonzero digits in a computation code.  Exponentiation lets us test it
by extracting one digit from a binomial expansion.
Large-base extraction of binomial coefficients goes back to Julia
Robinson \cite[Section~5]{robinson1952}.  Its combination with binomial
parity gives the binary-masking method used by
Jones and Matiyasevich \cite[Section~3]{jm1991}; the parity criterion is
the case modulo two of Lucas's theorem \cite[Section~3, p.~52]{lucas1878}.
We give its short proof and
retain all witnesses of the three power calls and two divisions.

\begin{proposition}\label{prop:subset}
The relation $x\bits y$ on $\NN^2$ has a complete single-fold
Diophantine representation.
\end{proposition}

\begin{proof}
Require $x\le y$ and use three calls to
Proposition~\ref{prop:power} to specify
\begin{equation}\label{eq:subset-powers}
 B=2^{y+1},\qquad H=B^x,\qquad N=(B+1)^y.
\end{equation}
All three bases are at least two.
Introduce natural variables $s,d,r,t$ and impose
\begin{equation}\label{eq:subset-digit}
 N=sBH+dH+r,\qquad d<B,\qquad r<H,\qquad d=2t+1.
\end{equation}
For fixed $B,H,N$, the equation and the two bounds determine
$s,d,r$ uniquely.  Indeed, division of $N$ by $H$ first determines
$r$ and the quotient $sB+d$; division of this quotient by $B$ then
determines $s,d$.

The binomial theorem gives
\[
 N=(B+1)^y=\sum_{j=0}^{y}\binom yj B^j.
\]
Every coefficient is at most $\sum_{j=0}^{y}\binom yj=2^y<B$,
so this is the base-$B$ expansion of $N$ without carries.
As $H=B^x$ and $0\le x\le y$, the digit $d$ extracted in
\eqref{eq:subset-digit} is therefore $\binom yx$.

To determine its parity, work in the polynomial ring over the field
$\mathbb F_2$ with two elements.  Repeated squaring gives
$(1+Z)^{2^i}=1+Z^{2^i}$, and hence
\[
 (1+Z)^y=\prod_{i:\,y_i=1}(1+Z^{2^i})
 \qquad\hbox{in }\mathbb F_2[Z].
\]
Each monomial in the product is obtained by choosing some of the
positions where $y_i=1$.  Uniqueness of binary expansion ensures
that distinct choices have distinct exponents.  The coefficient of
$Z^x$ is thus one precisely when $x\bits y$.
Equivalently, $\binom yx$ is odd precisely in that case.
The final equation $d=2t+1$ in \eqref{eq:subset-digit} proves the
required equivalence.

For a true instance, the three outputs in
\eqref{eq:subset-powers} and all their auxiliary variables are
unique by Proposition~\ref{prop:power}.  The two divisions uniquely
determine $s,d,r$, and oddness uniquely determines the natural number
$t=(d-1)/2$.  Every inequality slack is likewise unique.
Conversely, when the relation is false, either $x>y$ or the extracted
digit is even, so the system has no solution.
The case $x=y=0$ causes no exception: it gives
\[
 B=2,\quad H=N=1,\quad s=r=t=0,\quad d=1.
\]
Thus existence and uniqueness include every natural input pair.
\end{proof}
